\subsection{Sample Extraction from First-Order Logic Worlds}

We have observed that in certain situations the log-likelihood of a \firstOrderLogic{} world with respect to a \HybridFOLNetwork{} coincides with the likelihood of a data set in a propositional \HybridLogicNetwork{}.
Let us now investigate the extraction process of these data set.
%The decomposition of the likelihood suggests the following approach to generate samples from groundings:
%%We propose the following approach to generate datacores from groundings:
%\begin{itemize}
%    \item Define an importance formula $\impformula$, which we decompose in the basis CP decomposition and interpret each slice as the one-hot encoding of the datapoint.
%    \item Define for $\atomenumeratorin$ extraction formulas $\extformulaof{\atomenumerator}$ generating the atoms $\catvariableof{\atomenumerator}$.
%    %Predicates along with assignment of variables / constants to its positions.
%    \item Contract the groundings of each formula $\extformulaof{\atomenumerator}$ with the grounding of $\impformula$ to build a data core.
%\end{itemize}
%\subsubsection{Representation by Tensor Networks}
We model the extraction process as a relation between a tuple of individuals and the extracted world in the factored system of atoms $\catvariableof{\atomenumerator}$.

\begin{definition}
    \label{def:extractionRelation}
    Given a \firstOrderLogic{} world $\worldindices$, an importance formula $\impformula$ and extraction formulas $\extformulaof{\catenumerator}$ for $\catenumeratorin$, we define the extraction relation
    \begin{align*}
        \extractionrelation \subset \left(\symindstates\right) \otimes \left(\atomstates\right)
    \end{align*}
    by
    \begin{align*}
        \extractionrelation
        = \{ (\shortindindices, \shortcatindices)
        \wcols  \groundingofat{\impformula}{\indexedshortindvariables} = 1 \ncond \uniquantwrtof{\catenumeratorin}{\catindexof{\atomenumerator} = \extformulaofat{\atomenumerator}{\indexedshortindvariables}} \} \, .
    \end{align*}
\end{definition}

The encoding of an extraction relation is the tensor
\begin{align*}
    \bencodingofat{\extractionrelation}{\shortindvariables,\shortcatvariables} \subset \left(\indspace\right) \otimes \left(\atomspace\right) \,
\end{align*}
and drawn in a contraction diagram by
\begin{center}
    \begin{tikzpicture}[scale=0.35, yscale=1, thick] % , baseline = -3.5pt


    \draw[] (2,-1) -- (2,1) node[midway,left] {\colorlabelsize $\catvariableof{0}$};
    \node[anchor=center] (text) at (4,0) {$\cdots$};
    \draw[] (6,-1) -- (6,1) node[midway,right] {\colorlabelsize $\catvariableof{\atomorder-1}$};

    \draw (1,-1) rectangle (7,-3);
    \node[anchor=center] (text) at (4,-2) {$\bencodingof{\extractionrelation}$};

    \draw[] (2,-3) -- (2,-5) node[midway,left] {\colorlabelsize $\indvariableof{0}$};
    \node[anchor=center] (text) at (4,-4) {$\cdots$};
    \draw[] (6,-3) -- (6,-5) node[midway,right] {\colorlabelsize $\indvariableof{\indorder-1}$};


    \node[anchor=center] (text) at (10,-2) {${=}$};


    \begin{scope}
        [shift={(12,0)}]

        \draw[->-] (2.5,1) -- (2.5,3) node[midway,right] {\colorlabelsize $\catvariableof{0}$};
        \draw (1,-1) rectangle (4,1);
        \node[anchor=center] (text) at (2.5,0) {$\bencodingof{\groundingof{\extformulaof{0}}}$};
        \node[anchor=center] (text) at (2.5,-2) {$\cdots$};

        \node[anchor=center] (text) at (6.5,0) {$\cdots$};

        \draw[->-] (10.5,1) -- (10.5,3) node[midway,right] {\colorlabelsize $\catvariableof{\atomorder-1}$};
        \draw (8.75,-1) rectangle (12.25,1);
        \node[anchor=center] (text) at (10.5,0) {$\bencodingof{\groundingof{\extformulaof{\atomorder\shortminus1}}}$};
        \node[anchor=center] (text) at (10.5,-2) {$\cdots$};

        \draw[-<-] (13,-3) -- (3.5,-3) ;
        \draw[-<-] (13,-5) -- (1.5,-5) ;

        \drawvariabledot{11.5}{-3}
        \draw[->-] (11.5,-3) -- (11.5,-1);

        \drawvariabledot{9.5}{-5}
        \draw[->-] (9.5,-5) -- (9.5,-1);

        \drawvariabledot{3.5}{-3}
        \draw[->-] (3.5,-3) -- (3.5,-1);

        \drawvariabledot{1.5}{-5}
        \draw[->-] (1.5,-5) -- (1.5,-1);

        \drawvariabledot{7.5}{-3}
        \draw[-<-] (7.5,-3) -- (7.5,-7) node[right] {\colorlabelsize $\indvariableof{\indorder-1}$} ;

        \node[anchor=center] (text) at (6.5,-6) {$\cdots$};

        \drawvariabledot{5.5}{-5}
        \draw[-<-] (5.5,-5) -- (5.5,-7) node[left] {\colorlabelsize $\indvariableof{0}$} ;


        \draw (13,-2) rectangle (17,-6);
        \node[anchor=center] (text) at (15,-4) {$\bencodingof{\groundingof{\impformula}}$};
        \node[anchor=center] (text) at (12,-3.75) {$\vdots$};
        \draw[->-] (17,-4) -- (18,-4);
        \drawvariabledot{18}{-4}
        \draw[] (18,-4) -- (19,-4);
        \draw (19,-3) rectangle (21,-5);
        \node[anchor=center] (text) at (20,-4) {$\tbasis$};

    \end{scope}


\end{tikzpicture}
\end{center}
Here the contraction of $\bencodingof{\impformula}$ with the truth vector $\tbasis$ represents the matching condition posed by $\impformula$ when extracting pairs of individuals.

%% Empirical Distribution
The empirical distribution of the extracted data is then the normalized contraction leaving only the legs to the extracted atomic formulas open, that is
\begin{align*}
    \empdistribution
    = \frac{
        \contractionof{\bencodingof{\extractionrelation}}{\shortcatvariables}
    }{
        \contraction{\bencodingof{\extractionrelation}}
    }  \, .
\end{align*}
Here the number of extracted data is the denominator
\begin{align*}
    \datanum
    = \contraction{\bencodingof{\extractionrelation}}
    = \contraction{\bencodingofat{\impformula}{\headvariableof{\impformula},\shortindvariables},\tbasisat{\headvariableof{\impformula}}}\, .
\end{align*}

We depict this by
\begin{center}
    
\begin{tikzpicture}[scale=0.35, yscale=1, thick] % , baseline = -3.5pt


    \draw[->-] (2,-1) -- (2,1) node[midway,left] {\colorlabelsize $\catvariableof{0}$};
    \node[anchor=center] (text) at (4,0) {$\cdots$};
    \draw[->-] (6,-1) -- (6,1) node[midway,right] {\colorlabelsize $\catvariableof{\atomorder-1}$};

    \draw (1,-1) rectangle (7,-3);
    \node[anchor=center] (text) at (4,-2) {$\empdistribution$};
    \node[anchor=center] (text) at (-1,-2) {$\datanum \,\, \cdot $};

    \node[anchor=center] (text) at (10,-2) {${=}$};

    \begin{scope}
        [shift={(12,0)}]

        \draw[->-] (2.5,1) -- (2.5,3) node[midway,right] {\colorlabelsize $\catvariableof{0}$};
        \draw (1,-1) rectangle (4,1);
        \node[anchor=center] (text) at (2.5,0) {$\bencodingof{\groundingof{\extformulaof{0}}}$};
        \node[anchor=center] (text) at (2.5,-2) {$\cdots$};

        \node[anchor=center] (text) at (6.5,0) {$\cdots$};

        \draw[->-] (10.5,1) -- (10.5,3) node[midway,right] {\colorlabelsize $\catvariableof{\atomorder-1}$};
        \draw (8.75,-1) rectangle (12.25,1);
        \node[anchor=center] (text) at (10.5,0) {$\bencodingof{\groundingof{\extformulaof{\atomorder\shortminus1}}}$};
        \node[anchor=center] (text) at (10.5,-2) {$\cdots$};

        \draw[-<-] (13,-3) -- (3.5,-3) ;
        \draw[-<-] (13,-5) -- (1.5,-5) ;

        \drawvariabledot{11.5}{-3}
        \draw[->-] (11.5,-3) -- (11.5,-1);

        \drawvariabledot{9.5}{-5}
        \draw[->-] (9.5,-5) -- (9.5,-1);

        \drawvariabledot{3.5}{-3}
        \draw[->-] (3.5,-3) -- (3.5,-1);

        \drawvariabledot{1.5}{-5}
        \draw[->-] (1.5,-5) -- (1.5,-1);

        \draw (13,-2) rectangle (17,-6);
        \node[anchor=center] (text) at (15,-4) {$\bencodingof{\groundingof{\impformula}}$};
        \node[anchor=center] (text) at (12,-3.75) {$\vdots$};
        \draw[->-] (17,-4) -- (18,-4);
        \drawvariabledot{18}{-4}
        \draw[] (18,-4) -- (19,-4);
        \draw (19,-3) rectangle (21,-5);
        \node[anchor=center] (text) at (20,-4) {$\tbasis$};

    \end{scope}


\end{tikzpicture}
\end{center}

%\subsubsection{Decomposition of extracted data}

To connect with the empirical distribution introduced in \secref{sec:empDistribution} we now show how the empirical distribution extracted from the interpretations of the formulas $\impformula,\extformulas$ on a \firstOrderLogic{} world $\worldindices$ can be represented by tensor networks.

First of all, we decompose the importance formula into a basis $\cpformat$ format (see \charef{cha:sparseRepresentation}), that is a decomposition
\begin{align*}
    \groundingofat{\impformula}{\shortindvariables}
    = \contractionof{
        \{\legcoreofat{\indenumerator}{\indvariableof{\indenumerator},\datvariable} \wcols \indenumeratorin \}
    }{\shortindvariables}
\end{align*}
such that all $\legcoreofat{\indenumerator}{\indvariableof{\indenumerator},\datvariable}$ are directed and boolean tensors.
Here an auxiliary variables $\datvariable$ taking values in $[\datanum]$ is introduced, which we call the data variable, which enumerates the non-vanishing coordinates of $\groundingof{\impformula}$.
With this decomposition, we can understand the decomposition of $\groundingofat{\impformula}{\shortindvariables}$ as a basis encoding of an term selection map $\secdatamap$ with coordinate maps defined such that
\begin{align*}
    \bencodingofat{\secdatamap_{\indenumerator}}{\indvariableof{\indenumerator},\datvariable}
    = \legcoreofat{\indenumerator}{\indvariableof{\indenumerator},\datvariable} \, .
\end{align*}
We depict this decomposition by:
\begin{center}
    \begin{tikzpicture}[scale=0.35, thick] % , baseline = -3.5pt

    \begin{scope}
        [shift={(0,2)}]
        \draw[] (0,1)--(0,-1) node[midway,left] {\colorlabelsize $\indvariableof{0}$};
        \draw[] (1.5,1)--(1.5,-1) node[midway,left] {\colorlabelsize $\indvariableof{1}$};
        \node[anchor=center] (text) at (3,0) {$\cdots$};
        \draw[] (4,1)--(4,-1) node[midway,right] {\colorlabelsize $\indvariableof{\indorder\shortminus1}$};
    \end{scope}


    \draw (-1,1) rectangle (5,-1);
    \node[anchor=center] (text) at (2,0) {\corelabelsize ${\groundingof{\impformula}}$};


    \node[anchor=center] (text) at (7,0) {${=}$};


    \begin{scope}
        [shift={(10,2)}]



        \coordinate (conposseldec) at (4.5,-5.5);
        \drawvariabledot{4.5}{-5.5}

        \draw (conposseldec) -- (4.5,-7.5) node[midway, right] {\colorlabelsize ${\datvariable}$}; % Unclear, whether this is the best notation!
        \draw (3.5,-7.5) rectangle (5.5, -9.5);
        \node[anchor=center] (text) at (4.5,-8.5) {\corelabelsize $\ones$};

        \draw[-<-] (0,1) -- (0,-1) node[midway,left] {\colorlabelsize $\indvariableof{0}$};
        \draw (-1,-1) rectangle (1, -3);
        \node[anchor=center] (text) at (0,-2) {\corelabelsize $\secdatacoreof{0}$};
        \draw[-<-] (0,-3) to[bend right=20] (conposseldec);


        \draw[-<-] (3,1) -- (3,-1) node[midway,left] {\colorlabelsize $\indvariableof{1}$};
        \draw (2,-1) rectangle (4, -3);
        \node[anchor=center] (text) at (3,-2) {\corelabelsize $\secdatacoreof{1}$};
        \draw[-<-] (3,-3) to[bend right=20]  (conposseldec);

        \node[anchor=center] (text) at (6,-2) {$\cdots$};

        \draw[-<-] (9,1) -- (9,-1) node[midway,left] {\colorlabelsize $\indvariableof{\indorder-1}$};
        \draw (7.75,-1) rectangle (10.25, -3);
        \node[anchor=center] (text) at (9,-2) {\corelabelsize $\secdatacoreof{\indorder-1}$};
        \draw[-<-] (9,-3) to[bend left=20]  (conposseldec);

    \end{scope}


\end{tikzpicture}
\end{center}

Based on these construction, we now provide a tensor network decomposition of the extracted empirical distribution.

\begin{theorem}
    \label{the:extractionrelationDecomposition}
    Given a \firstOrderLogic{} world $\worldindices$, an importance formula $\impformula$ and extraction formulas $\extformulaof{\catenumerator}$ for $\catenumeratorin$, we have
    \begin{align*}
        \bencodingofat{\extractionrelation}{\shortindvariables,\shortcatvariables} =
        \contractionof{
            \{\bencodingofat{\groundingof{\extformulaof{\atomenumerator}}}{\catvariableof{\catenumerator},\shortindvariables} \wcols \catenumeratorin\}
            \cup \{\bencodingofat{\secdatamap_{\indenumerator}}{\indvariableof{\indenumerator},\datvariable} \wcols \indenumeratorin\}
        }{\shortindvariables,\shortcatvariables}
    \end{align*}
    and thus
    \begin{align*}
        \empdistributionat{\shortcatvariables} =
        \frac{1}{\datanum}  \contractionof{
            \{\bencodingofat{\groundingof{\extformulaof{\atomenumerator}}}{\catvariableof{\catenumerator},\shortindvariables} \wcols \catenumeratorin\}
            \cup \{\bencodingofat{\secdatamap_{\indenumerator}}{\indvariableof{\indenumerator},\datvariable} \wcols \indenumeratorin\}
        }{\shortcatvariables} \, .
    \end{align*}
\end{theorem}
\begin{proof}
    To show the first claim, let us choose arbitrary state tuples $\shortindindices$ and $\shortcatindices$.
    We then have
    \begin{align*}
        &\contractionof{
            \{\bencodingofat{\groundingof{\extformulaof{\atomenumerator}}}{\catvariableof{\catenumerator},\shortindvariables} \wcols \catenumeratorin\}
            \cup \{\bencodingofat{\secdatamap_{\indenumerator}}{\indvariableof{\indenumerator},\datvariable} \wcols \indenumeratorin\}
        }{\indexedshortindvariables,\indexedshortcatvariables} \\
        & \quad  =  \contraction{
            \{\bencodingofat{\groundingof{\extformulaof{\atomenumerator}}}{\indexedcatvariableof{\catenumerator},\indexedshortindvariables} \wcols \catenumeratorin\}
            \cup \{\bencodingofat{\secdatamap_{\indenumerator}}{\indexedindvariableof{\indenumerator},\datvariable} \wcols \indenumeratorin\}
        } \, .
    \end{align*}
    This contraction evaluates to $1$, if and only if for all $\catenumeratorin$ we have $\bencodingofat{\groundingof{\extformulaof{\atomenumerator}}}{\catvariableof{\catenumerator},\shortindvariables}=1$ and
    \begin{align*}
        \contraction{\{\bencodingofat{\secdatamap_{\indenumerator}}{\indexedindvariableof{\indenumerator},\datvariable} \wcols \indenumeratorin\}}  = 1 \, .
    \end{align*}
    The first condition is equal to $\catindexof{\atomenumerator} = \extformulaofat{\atomenumerator}{\indexedshortindvariables}$ for all $\catenumeratorin$ and the second to
    \begin{align*}
        \groundingofat{\impformula}{\indexedshortindvariables} = 1 \, .
    \end{align*}
    Comparing with the definition of the extraction relation (see \defref{def:extractionRelation}), we notice that these conditions are equal to $(\shortindindices,\shortcatindices)\in\extractionrelation$ and therefore to
    \begin{align*}
        \bencodingofat{\extractionrelation}{\indexedshortindvariables,\indexedshortcatvariables} \, .
    \end{align*}
    The first claim follows, since $\bencodingof{\extractionrelation}$ is boolean, as is the contraction of the cores $\bencodingof{\groundingof{\extformulaof{\atomenumerator}}}$ with the cores $\bencodingof{\secdatamap_{\indenumerator}}$, which leaves the outgoing variables $\shortcatvariables$ open.
    The second claim follows from the first using that $\empdistributionat{\shortcatvariables}=\frac{1}{\datanum}\contractionof{\bencodingof{\extractionrelation}}{\shortcatvariables}$.
\end{proof}

To connect with the representation of empirical distributions based on data cores (see \secref{sec:empDistribution}), we now form data cores by contractions with the grounding of extraction formulas with the cores $\bencodingof{\secdatamap_{\indenumerator}}$ (see \figref{fig:datacoreGeneration}),
\begin{align*}
    \datacoreofat{\atomenumerator}{\catvariableof{\catenumerator},\datvariable}
    = \contractionof{
        \{\bencodingofat{\groundingof{\extformulaof{\atomenumerator}}}{\catvariableof{\catenumerator},\shortindvariables}\}
        \cup \{ \legcoreofat{\indenumerator}{\indvariableof{\indenumerator},\datvariable} \wcols \indenumeratorin\}
    }{\catvariableof{\atomenumerator},\datvariable} \, .
\end{align*}

% Empirical distribution
The empirical distribution is then a tensor network of these tensors, as we show next.

\begin{theorem}
    \label{the:extractionDataCores}
    We have
    \begin{align*}
        \contractionof{\bencodingof{\extractionrelation}}{\shortcatvariables}
        = \contractionof{\{\datacoreofat{\atomenumerator}{\datvariable,\catvariableof{\atomenumerator}} \wcols \atomenumeratorin\}}{\shortcatvariables}
    \end{align*}
    and thus
    \begin{align*}
        \empdistributionat{\shortcatvariables}
        = \frac{1}{\datanum} \contractionof{\{\datacoreofat{\atomenumerator}{\datvariable,\catvariableof{\atomenumerator}}  \wcols \atomenumeratorin\}}{\shortcatvariables} \, .
    \end{align*}
\end{theorem}
\begin{proof}
    By \theref{the:extractionrelationDecomposition} we have
    \begin{align*}
        \bencodingofat{\extractionrelation}{\shortindvariables,\shortcatvariables} =
        \contractionof{
            \{\bencodingofat{\groundingof{\extformulaof{\atomenumerator}}}{\catvariableof{\catenumerator},\shortindvariables} \wcols \catenumeratorin\}
            \cup \{\bencodingofat{\secdatamap_{\indenumerator}}{\indvariableof{\indenumerator},\datvariable} \wcols \indenumeratorin\}
        }{\shortindvariables,\shortcatvariables} \, .
    \end{align*}
    Since $\bencodingofat{\secdatamap_{\indenumerator}}{\indvariableof{\indenumerator},\datvariable}$ are directed and boolean, they can be copied and separately contracted with each $\groundingof{\extformulaof{\atomenumerator}}$, without changing the contraction.
    We arrive at
    \begin{align*}
        &\bencodingofat{\extractionrelation}{\shortindvariables,\shortcatvariables} \\
        &\quad = \contractionof{
            \big\{\contractionof{
                \{\bencodingofat{\groundingof{\extformulaof{\atomenumerator}}}{\catvariableof{\catenumerator},\shortindvariables}\}
                \cup \{\bencodingofat{\secdatamap_{\indenumerator}}{\indvariableof{\indenumerator},\datvariable} \wcols \indenumeratorin\}
            }{\catvariableof{\catenumerator},\datvariable} \wcols \catenumeratorin \big\}
        }{\shortindvariables,\shortcatvariables} \\
        & \quad =  \contractionof{\{\datacoreofat{\atomenumerator}{\datvariable,\catvariableof{\atomenumerator}}  \wcols \atomenumeratorin\}}{\shortcatvariables} \, ,
    \end{align*}
    which established the claim.
\end{proof}

% Efficient contraction: Do also basis decomposition of the extraction query and use efficient contraction!
%Towards efficient calculation of the data cores, we build a basis CP decomposition of $\groundingof{\impformula}$, where we further demand $\scalarcore=\ones$.
%This is a collection of basis leg cores $\legcoreof{\fixedimpformula,\indenumerator}$ such that
%\begin{align*}
%    \fixedimpformula[\shortindvariablelist]
%    = \contractionof{ \left\{ \legcoreofat{\fixedimpformula,\indenumerator}{\datvariable,\indvariableof{\indenumerator}} \wcols \indenumeratorin \right\} }{\shortindvariablelist} \, .
%\end{align*}

% Data enumeration -> To representation
%We can further utilize any decomposition of $\impformula$ into a directed and binary CP Format to enumerate the datapoints by the slice index $\datindex$. % Approaches like SPARQL directly give us these by solution mappings.
%Understanding $\impformula$ as a query on the world being the database, such decomposition is given by the set of solution mappings.


\begin{figure}[t]
    \begin{center}
        \begin{tikzpicture}[scale=0.35, yscale=1, thick] % , baseline = -3.5pt


    \draw[->-] (4,-1) -- (4,1) node[midway, right] {\colorlabelsize $\catvariableof{\atomenumerator}$};
    \draw (3,-1) rectangle (5,-3);
    \node[anchor=center] (text) at (4,-2) {$\datacoreof{\atomenumerator}$};
    \draw[-<-] (4,-3) -- (4,-5) node[midway, right] {\colorlabelsize $\datvariable$};

    \node[anchor=center] (text) at (7,-2) {${=}$};

    \begin{scope}
        [shift={(10,0)}]

        \draw[->-] (3,1) -- (3,3) node[midway, right] {\colorlabelsize $\catvariableof{\atomenumerator}$};
        \draw (-1,1) rectangle (7,-1);
        \node[anchor=center] (text) at (3,0) {$\bencodingof{\groundingof{\extformulaof{\atomenumerator}}}$};

        \draw[->-] (0,-3) -- (0,-1) node[midway,left] {\colorlabelsize $\indvariableof{0}$};
        \draw[->-] (3,-3) -- (3,-1) node[midway,left] {\colorlabelsize $\indvariableof{1}$};
        \draw[->-] (6,-3) -- (6,-1) node[midway,left] {\colorlabelsize $\indvariableof{2}$};


    \end{scope}

    \begin{scope}
        [shift={(10,-2)}]

        \coordinate (conposseldec) at (4.5,-5.5);
        \drawvariabledot{4.5}{-5.5}
        \draw[-<-] (conposseldec) -- (4.5,-7.5) node[midway, right] {\colorlabelsize $\indexvariable$};

        \draw (-1,-1) rectangle (1, -3);
        \node[anchor=center] (text) at (0,-2) {\corelabelsize $\bencodingof{\secdatamap_0}$};%{\corelabelsize $\legcoreof{\fixedimpformula,0}$};
        \draw[-<-] (0,-3) to[bend right=20] (conposseldec);

        \draw (2,-1) rectangle (4, -3);
        \node[anchor=center] (text) at (3,-2) {\corelabelsize $\bencodingof{\secdatamap_1}$};%{\corelabelsize $\legcoreof{\fixedimpformula,1}$};
        \draw[-<-] (3,-3) to[bend right=20]  (conposseldec);

        \draw (5,-1) rectangle (7, -3);
        \node[anchor=center] (text) at (6,-2) {\corelabelsize $\bencodingof{\secdatamap_2}$};%{\corelabelsize $\legcoreof{\fixedimpformula,2}$};
        \draw[-<-] (6,-3) to[bend right=-20]  (conposseldec);

        \draw[-<-] (9,1) -- (9,-1) node[midway,left] {\colorlabelsize $\indvariableof{3}$};
        \draw (8,-1) rectangle (10, -3);
        \node[anchor=center] (text) at (9,-2) {\corelabelsize $\bencodingof{\secdatamap_3}$};%{\corelabelsize $\legcoreof{\fixedimpformula,3}$};
        \draw[-<-] (9,-3) to[bend right=-20]  (conposseldec);


        \node[anchor=center] (text) at (12,-2) {$\cdots$};

        \draw[-<-] (15,1) -- (15,-1) node[midway,left] {\colorlabelsize $\indvariableof{\indorder-1}$};
        \draw (13.5,-1) rectangle (16.5, -3);
        \node[anchor=center] (text) at (15,-2) {\corelabelsize $\bencodingof{\secdatamap_{\indorder-1}}$};%{\corelabelsize $\legcoreof{\fixedimpformula,\variableorder-1}$};
        \draw[-<-] (15,-3) to[bend left=20]  (conposseldec);


        \draw (8,1) rectangle (16, 3);
        \node[anchor=center] (text) at (12,2) {\corelabelsize $\ones$};


    \end{scope}


    \node[anchor=center] (text) at (29,-2) {${=}$};


    \begin{scope}
        [shift={(32,0)}]

        \draw[->-] (3,1) -- (3,3) node[midway, right] {\colorlabelsize $\catvariableof{\atomenumerator}$};
        \draw (-1,1) rectangle (7,-1);
        \node[anchor=center] (text) at (3,0) {$\bencodingof{\groundingof{\extformulaof{\atomenumerator}}}$};

        \draw[->-] (0,-3) -- (0,-1) node[midway,left] {\colorlabelsize $\indvariableof{0}$};
        \draw[->-] (3,-3) -- (3,-1) node[midway,left] {\colorlabelsize $\indvariableof{1}$};
        \draw[->-] (6,-3) -- (6,-1) node[midway,left] {\colorlabelsize $\indvariableof{2}$};


    \end{scope}

    \begin{scope}
        [shift={(32,-2)}]


        \coordinate (conposseldec) at (3,-5.5);
        \drawvariabledot{3}{-5.5}
        \draw[-<-] (conposseldec) -- (3,-7.5) node[midway, right] {\colorlabelsize $\datvariable$};

        \draw (-1,-1) rectangle (1, -3);
        \node[anchor=center] (text) at (0,-2){\corelabelsize $\bencodingof{\secdatamap_0}$};%{\corelabelsize $\legcoreof{\fixedimpformula,0}$};
        \draw[-<-] (0,-3) to[bend right=20] (conposseldec);

        \draw (2,-1) rectangle (4, -3);
        \node[anchor=center] (text) at (3,-2) {\corelabelsize $\bencodingof{\secdatamap_1}$};%{\corelabelsize $\legcoreof{\fixedimpformula,1}$};
        \draw[-<-] (3,-3) to[bend right=0]  (conposseldec);

        \draw (5,-1) rectangle (7, -3);
        \node[anchor=center] (text) at (6,-2) {\corelabelsize $\bencodingof{\secdatamap_2}$};%{\corelabelsize $\legcoreof{\fixedimpformula,2}$};
        \draw[-<-] (6,-3) to[bend right=-20]  (conposseldec);


    \end{scope}


\end{tikzpicture}
    \end{center}
    \caption{Generation of a data core for the variable $\catvariableof{\catenumerator}$ given an extraction formula $\extformulaof{\catenumerator}$ and an importance formula, which grounding is decomposed into a basis CP format with leg vectors $\bencodingofat{\secdatamap_{\indenumerator}}{\indvariableof{\indenumerator},\datvariable}$.
    Term variables, which are appearing in the importance formula, but not in the extraction formula $\extformulaof{\catenumerator}$ can be treated trivally by contraction with the trivial tensor (here $\indvariableof{4},\ldots,\indvariableof{\indorder-1})$.
    }
    \label{fig:datacoreGeneration}
\end{figure}


% Comment: Exploitation of common structure
When many atom extraction formulas differ only by a constant, we can replace the constant by an auxiliary term variable.
The atoms are then the atomizations of this variable (see \secref{sec:categoricalTN}), treated as a categorical variable, with respect to the constant in the extraction query.
The advantages are that we can avoid the $\bencodingof{}$-formalism and directly model the categorical distributions.

This also enables a batchwise computation of multiple $\sparql$ queries, which differ only in one constant.


%\subsection{Design of the Formulas}
%
%Most intuitive when labeling individuals by classes.
%Extraction formulas $\extformulas$ can then be defined by subclasses of the member of a class and relations between objects of different classes. % Koller calls atomic formulas the template attributes
%We then choose $\formulaset$ as more involved formulas decomposed into connectives acting on these atoms.
%The importance formula $\impformula$ is then designed based on class memberships to ensure, that the arguments of the formulas are always of specific classes. % Koller specifies to each argument of the attributes a class
%
%% Approach
%We propose to
%\begin{itemize}
%    \item Execute an extraction query to get pairs of individuals (the pairDf).
%    \item Propositionalize the FOL Formulas independently on each tuple taking the individuals as a set of constant and filtering on the possible properties of each individuals.
%    (Can understand as adding knowledge that most of the relations do not hold)
%    \item Understand each such generated knowledge base as datapoint and average over them to get the empirical distribution to be fit.
%    \item Fit a MLN describing the statistical relations of unseen results of the extraction query, based on likelihood maximation.
%\end{itemize}